% ─────────────────────────────────────────────────────────────────────────────
\subsection{Embedded Metadata Sensitivity Check}
\label{subsec:results-embedded-metadata-sensitivity}
% ─────────────────────────────────────────────────────────────────────────────

After normalization of the selected black-box outputs, a secondary sensitivity
check was performed on the residual metadata embedded in the blind-view image
files. Bundle blindness was enforced through opaque file names, paths, and
directory structure. Embedded metadata was preserved, however, because some
software systems may index information stored within the image bytes.

The purpose of this analysis is limited to measuring how the reported aggregate
metrics change when files containing predefined semantically or experimentally
sensitive metadata terms are excluded. It does not establish whether any
software system accessed those fields or whether metadata causally influenced
an individual decision.

The audit covered all 11\,500 files in the forensic evaluation bundle. Embedded
metadata was detected in 7\,640 files, while 15 files contained at least one
predefined sensitive-term match. These 15 files comprise 11 clean binary images
and four \gls{ood} images. No adversarial or anti-forensic artifact contains a
sensitive-term match. Consequently, the identified metadata channel does not
directly affect the perturbation-family populations used for the principal
robustness comparisons.

Table~\ref{tab:results-embedded-metadata-scope} summarizes the affected
populations and the corresponding leave-out denominators.

\begin{table}[H]
\centering
\scriptsize
\renewcommand{\arraystretch}{1.15}
\setlength{\tabcolsep}{4pt}
\caption{Scope of the residual embedded-metadata sensitivity check}
\label{tab:results-embedded-metadata-scope}
\begin{tabularx}{\textwidth}{@{}p{0.25\textwidth} r r r X@{}}
\toprule
\textbf{Population} &
\textbf{Flagged files} &
\textbf{Original denominator} &
\textbf{Leave-out denominator} &
\textbf{Interpretive scope} \\
\midrule

Complete bundle
& 15
& 11\,500
& 11\,485
& Describes the prevalence of predefined metadata-term matches in the complete
blind-view bundle. \\

Clean binary subset
& 11
& 1\,000
& 989
& Supports a leave-out comparison of clean accuracy for each black-box
configuration. \\

Adversarial family
& 0
& 5\,000
& 5\,000
& No sensitive-term match occurs in the adversarial or adversarial-style
artifact populations. \\

Anti-forensic family
& 0
& 5\,000
& 5\,000
& No sensitive-term match occurs in the anti-forensic artifact populations. \\

\gls{ood} branch
& 4
& 500
& 496
& Supports a leave-out comparison of the observable \gls{ood} weapon assignment
rate. \\

\bottomrule
\end{tabularx}
\end{table}

Table~\ref{tab:results-embedded-metadata-leaveout} reports the exact metric
changes after excluding the affected files. The clean-accuracy delta is defined
as the leave-out value minus the original value. The \gls{ood} delta follows the
same convention for the weapon assignment rate.

\begin{table}[H]
\centering
\scriptsize
\renewcommand{\arraystretch}{1.12}
\setlength{\tabcolsep}{3.2pt}
\caption{Leave-out sensitivity results for files containing predefined
embedded-metadata term matches}
\label{tab:results-embedded-metadata-leaveout}
\resizebox{\textwidth}{!}{%
\begin{tabular}{@{}l r r r r r r@{}}
\toprule
\textbf{Configuration} &
\textbf{Clean original} &
\textbf{Clean leave-out} &
\textbf{Clean delta} &
\textbf{\acrshort{ood} original} &
\textbf{\acrshort{ood} leave-out} &
\textbf{\acrshort{ood} delta} \\
\midrule

\gls{magnetaxiomtool}/\gls{magnetai}
& 0.9500 & 0.9515 & +0.0015
& 0.3600 & 0.3548 & -0.0052 \\

\gls{excirefoto2025} \texttt{D20}
& 0.9270 & 0.9282 & +0.0012
& 0.2380 & 0.2379 & -0.0001 \\

\gls{excirefoto2025} \texttt{D50}
& 0.9440 & 0.9454 & +0.0014
& 0.3400 & 0.3347 & -0.0053 \\

\gls{excirefoto2025} \texttt{D80}
& 0.9150 & 0.9141 & -0.0009
& 0.5220 & 0.5181 & -0.0039 \\

\gls{cellebriteinseyets}
& 0.9640 & 0.9636 & -0.0004
& 0.2920 & 0.2883 & -0.0037 \\

\gls{griffeye}/\gls{t3kcore}
& 0.9780 & 0.9788 & +0.0008
& 0.2600 & 0.2601 & +0.0001 \\

\bottomrule
\end{tabular}%
}

\vspace{0.4em}
\footnotesize
\textit{Note.} Clean values use 1\,000 original and 989 leave-out binary
inputs. \acrshort{ood} values use 500 original and 496 leave-out inputs. Deltas
are leave-out minus original. The analysis measures metric sensitivity only and
does not identify a causal contribution of embedded metadata to any black-box
decision.
\end{table}

The clean-accuracy changes range from \(-0.0009\) to \(+0.0015\). The largest
absolute change in the \gls{ood} weapon assignment rate is approximately
0.0053. These variations do not alter the qualitative ordering or the principal
operational conclusions reported for the selected configurations.

The result must nevertheless not be interpreted as proof that embedded metadata
was irrelevant to every affected decision. For the four \gls{ood} files with
sensitive-term matches, the observable outputs do not permit separation of
visual and metadata contributions. Similarly, the fact that some clean
\texttt{weapon} images remain unflagged despite weapon-related metadata shows
only that the presence of such a term is not sufficient to determine every
normalized output.

A causal assessment would require a paired metadata-neutralized bundle in which
the image pixels remain unchanged while the embedded fields are removed or
standardized. Such a paired experiment was not included in the frozen protocol.
The identified cases are therefore retained in the corpus and documented as a
residual limitation of bundle blindness rather than removed retrospectively.
